\documentclass{article}
\usepackage[utf8]{inputenc}
\usepackage{amsmath}
\usepackage{amsfonts}
\usepackage{amssymb}
\usepackage{geometry}
\usepackage{tcolorbox}
\usepackage{listings}
\usepackage{xcolor}
\usepackage{hyperref}

\geometry{a4paper, margin=1in}

\definecolor{codegreen}{rgb}{0,0.6,0}
\definecolor{codegray}{rgb}{0.5,0.5,0.5}
\definecolor{codepurple}{rgb}{0.58,0,0.82}
\definecolor{backcolour}{rgb}{0.95,0.95,0.92}

\lstset{
    backgroundcolor=\color{backcolour},   
    commentstyle=\color{codegreen},
    keywordstyle=\color{magenta},
    numberstyle=\tiny\color{codegray},
    stringstyle=\color{codepurple},
    basicstyle=\ttfamily\small,
    breakatwhitespace=false,         
    breaklines=true,                 
    captionpos=b,                    
    keepspaces=true,                 
    numbers=left,                    
    numbersep=5pt,                  
    showspaces=false,                
    showstringspaces=false,
    showtabs=false,                  
    tabsize=2
}

\title{Module 04: Feature Engineering II - Automated Selection and Adversarial Validation}
\author{Cybernauts 2026 Datathon Campaign}
\date{May 2026}

\begin{document}

\maketitle

\begin{abstract}
As you engineer hundreds of features, your dataset becomes cluttered with noise and redundancy. This note covers the critical "Pruning Phase"---how to systematically drop features that don't help, and a secret Grandmaster technique called "Adversarial Validation" to ensure your model doesn't fail when it hits the hidden test set.
\end{abstract}

\tableofcontents
\newpage

\section{Systematic Pruning: Cleaning the Feature Forest}
Engineering 200 features is easy; keeping the right 50 is the hard part.

\subsection{Feature Importance Filtering}
The fastest way to rank features is by running a baseline \textbf{LightGBM} model with default parameters.
\textit{Why?} Tree-based models calculate \textbf{Gain} and \textbf{Split} importance, telling you exactly which columns provided the most signal.

\begin{lstlisting}[language=Python]
import lightgbm as lgb
model = lgb.LGBMRegressor().fit(X, y)
importance_df = pd.DataFrame({'feature': X.columns, 'importance': model.feature_importances_})
importance_df = importance_df.sort_values(by='importance', ascending=False)
# Strategy: Drop the bottom 30% of features
\end{lstlisting}

\subsection{Dropping Collinear Features}
Redundant features add noise and slow down training. If two features have a correlation $> 0.95$, they are essentially the same.

\begin{tcolorbox}[colback=orange!5!white,colframe=orange!75!black,title=The Selection Rule]
Identify pairs of highly correlated features. Keep the one that has the \textbf{strongest relationship to the target} and drop the other.
\end{tcolorbox}

\section{Adversarial Validation: Detecting Data Drift}
This is a high-level competitive technique to see if your training data and test data come from different distributions.

\subsection{The Concept}
\begin{enumerate}
    \item Label all your training rows as \texttt{0} and all test rows as \texttt{1}.
    \item Combine them into one big dataset (dropping the actual target).
    \item Train a model to predict: \textbf{"Is this row from the Test Set?"}
\end{enumerate}

\subsection{Interpreting the Results}
\begin{itemize}
    \item \textbf{AUC $\approx$ 0.50:} The model can't tell the difference. Your data is stable.
    \item \textbf{AUC $>$ 0.70:} The model can identify test rows. This means some features behave differently in the test set.
\end{itemize}

\begin{tcolorbox}[colback=red!5!white,colframe=red!75!black,title=Action Plan]
If the AUC is high, check the feature importance of your Adversarial Model. The features at the top are "leaking" information about the test set. \textbf{Drop them immediately} to avoid a massive score drop on the private leaderboard.
\end{tcolorbox}

\section{Practice Problems}
\textbf{P1:} You have engineered 150 features. You run a baseline model and find that 40 features have an importance score of exactly 0. What should you do? \\
\textit{Answer: Drop them. They are pure noise and never used for a single split in the tree model.}

\textbf{P2:} Your Adversarial Validation model has an AUC of 0.85. The top feature is "Transaction\_Date". What does this tell you? \\
\textit{Answer: The test data likely comes from a different time period than the training data. This feature allows the model to "memorize" the time gap. You should probably drop it or transform it into a relative feature (e.g., "Days Since Start").}

\end{document}
